\documentclass[10pt,a4paper]{article}
\usepackage[T1]{fontenc}
\usepackage[margin=23mm]{geometry}
\usepackage{lmodern,mathrsfs,amsmath,amssymb,amsthm,mathtools,booktabs,microtype}
\usepackage[hidelinks]{hyperref}
\newtheorem{theorem}{Theorem}
\newtheorem{lemma}[theorem]{Lemma}
\emergencystretch=2em
\title{Split-zero inversion and a finite divisor-capacity theorem for ES}
\author{The Clankers}
\date{15 September 2026 --- research preprint}
\begin{document}\maketitle
\begin{abstract}
We distinguish sparse support-preserving inversion from ordinary amplitude
inversion in a split-zero coefficient algebra. Positive bounded-divisor factors
then give exact formulas for all prime multiplicities. A complete computer-assisted
classification at residual 27 proves that any occupied shell has a witness
using a reserved budget of at most nine prime occurrences in at most three
nonidentity residue classes. The three-class bound is necessary at a certified
hard prime. Source arithmetic, labels and inverse fibres remain explicit;
no universal ES or historical-priority assertion is made.
\end{abstract}

\paragraph{The scalar, not a substitute support index.}
For a nonzero commutative ring $A$, let
$G(A)=\{\tau\}\sqcup A^\bullet$, with $e=0_A^\bullet\ne\tau$.
This is the source's absent/supported distinction \cite{SZ}. A sparse vector
in $G(A)[\Gamma]$ is a pair $(S,f)$ with $S\subseteq\Gamma$ and
$\operatorname{nz}(f)\subseteq S$; a zero amplitude in $S$ is $e$.
Its operations are
\[
 (S,f)+(T,g)=(S\cup T,f+g),\qquad(S,f)(T,g)=(ST,fg).
 \tag{1}
\]
An amplitude sum may cancel; its presence does not. The inverse representation
reads $f$ on $S$ and $\tau$ off $S$.

\begin{theorem}[The precise inverse boundary]
The global units of $G(A)[\Gamma]$ are exactly $a^\bullet[g]$,
$a\in A^\times$. In the ordinary real group algebra, an element and its inverse
cannot both have nonnegative coefficients unless both are monomials.
\end{theorem}
\begin{proof}
If a product is the global identity, its support product is the singleton
$\{1\}$. Fixing a member of either nonempty support forces the other support
to be a singleton; the amplitudes must be units. Nonnegative real convolution
has the same noncancellation property. The converses are immediate.
\end{proof}
This is not a prohibition of linear inverses inside a fixed nonbottom fibre.
For example, in additive $C_6$ the ordinary inverse identity
\[
 ([3]+2[1])\frac{[3]-2[1]+4[5]}9=[0]
\]
returns instead $1^\bullet[0]+e[2]+e[4]=E_H$ in the sparse split algebra,
$H=\{0,2,4\}$. It is idempotent and acts as an identity only after the
specified coset saturation. In particular
$E_H([3]+2[1])=[3]+2[1]+e[5]$.
The earlier full Fourier support of this ordinary target \cite{Target} is
not disputed. The added present-zero coordinate is not a positive divisor count.

\paragraph{Actual arithmetic factors.}
Fix $R\equiv3\pmod4$, $a>R/2$, $\gcd(a,R)=1$, and $n=4a-R$.
Write $a=\prod q_i^{e_i}$, keeping the actual primes and their multiplicities.
In $\Gamma_R=(\mathbb Z/R\mathbb Z)^\times$ define
\[
 P^E_{g,e}=\sum_{f=0}^{2e}[g]^f,\qquad
 P^M_{g,e}=\sum_{\beta=-e}^{e}[g]^\beta.
 \tag{2}
\]
The E target is $-4^{-1}$, with $u=\prod q_i^{f_i}$; the M target is $-1$,
with $u=a\prod q_i^{\beta_i}$. Thus a target coefficient in
$\prod_iP^c_{q_i,e_i}$ counts exactly the original channel-$c$ states.
These positive count coefficients have a support-faithful sparse lift to
$G(\mathbb Q)$; no signed cancellation can create a spurious coefficient.

\begin{theorem}[All exponents from finite exact data]
In the ordinary formal group algebra,
\[
 \sum_{e\ge0}P^E_{g,e}z^e=\frac{1+z[g]}{(1-z)(1-z[g]^2)},\qquad
 \sum_{e\ge0}P^M_{g,e}z^e=\frac{1+z}{(1-z[g])(1-z[g]^{-1})}.
 \tag{3}
\]
Each has a specified positive expansion with one-to-one coefficient indices.
If $o=\operatorname{ord}(g)$ and $e=r+ok$, $0\le r<o$, then
\[
 P^c_{g,e}=P^c_{g,r}+2k\sum_{h\in\langle g\rangle}[h].
 \tag{4}
\]
Its support is unchanged by $e\mapsto\min(e,\lfloor o/2\rfloor)$.
\end{theorem}
\begin{proof}
The E expansion indices $(j,k,\varepsilon)\in\mathbb N^2\times\{0,1\}$
map to $(e,f)=(j+k+\varepsilon,2k+\varepsilon)$; the inverse is
$\varepsilon=f\bmod2$, $k=(f-\varepsilon)/2$, $j=e-k-\varepsilon$.
For M the map is $(e,\beta)=(j+k+\varepsilon,j-k)$, with inverse
$\varepsilon=(e-\beta)\bmod2$, $j=(e-\varepsilon+\beta)/2$,
$k=(e-\varepsilon-\beta)/2$. These are nonnegative exactly on the original
bounds. Increasing $e$ by $o$ adds two full cycles of $g$, proving (4).
A consecutive interval fills its cyclic subgroup when $2e+1\ge o$.
\end{proof}
The positive expansions, not formal subtraction inverses, are lifted to the
split carrier. Equation (4) is an exact multiaffine quasipolynomial formula
for arbitrary multiplicities. The cap preserves availability, not counts.

\paragraph{Grouping retains its complete fibres.}
For distinct primes of the same residue, aggregate capacity is
$E_g=\sum_{q\equiv g}e_q$. The support is the interval of total capacity;
its count polynomial is instead $\prod_{q\equiv g}(1+T+\cdots+T^{2e_q})$.
Every coefficient fibre is the set of integer compositions
$\sum f_q=f$, $0\le f_q\le2e_q$. It is nonempty for every $0\le f\le2E_g$,
by successive addition of integer intervals. These compositions supply all
actual prime-exponent lifts. For example, one prime squared has five choices,
whereas two distinct same-residue primes once each have nine.

\begin{theorem}[Complete residual-27 capacity classification]
Put $E_j(a)=\sum_{q\equiv2^j\ (27)}v_q(a)$ for $1\le j\le17$.
There is a fully specified list of 174 coordinatewise minimal profiles such
that a residual-27 E or M hit exists exactly when $(E_j(a))$ dominates a
listed profile. Nine profiles use one residue class, 99 use two, and 66 use
three; none uses more. Their total capacities are at most nine.
Every hit therefore has an original witness reserving at most nine actual
prime occurrences from at most three nonidentity residue classes.
\end{theorem}
\begin{proof}[Proof and complete finite certificate]
The group has order 18, generated by 2; the target logarithms are 7 and 9.
Group equal-residue factors using their composition fibres, then cap coordinate
$j$ at $c_j=\lfloor(18/\gcd(j,18))/2\rfloor$ using Theorem~2.
Occupancy increases coordinatewise.

For the $2^{17}$ zero/one supports, use 18-bit masks, initially $\{0\}$.
Adding class $j$ replaces E-mask $W$ by $W\cup(W+j)\cup(W+2j)$ and M-mask by
$W\cup(W+j)\cup(W-j)$, all modulo 18. Exactly 360 supports miss: 256 have only
even logarithms, which miss at every multiplicity. On the other 104 supports,
enumerate all positive capped coordinates; there are 14,040 profiles. Keep a
hit precisely when all its immediate predecessors miss. Including the 85
squarefree minimal profiles gives the 174 entries and the asserted bounds.
Every entry stores an actual target exponent word.

An independent original-residue checker starts at the zero capacity vector
and visits every capped unoccupied vector. It finds 62,812 such vectors and
tests 629,841 upward frontier edges. All 567,030 occupied boundary edges dominate
a verified listed profile. Every increasing path either remains in the
unoccupied set or first crosses such a boundary; hence no minimal hit is
omitted. Both exhaustive algorithms are included with their exact finite
bounds. This is the finite computer-assisted part of the theorem, not a
finite-prime extrapolation.

For a dominated profile reserve its capacities among actual occurrences in
the original factorization and allocate its target word using the composition
fibres. Unreserved E exponents are zero; unreserved M centered exponents are
zero. The original $a$ is unchanged. The reconstruction below completes the proof.
\end{proof}
The 104 odd-logarithm misses have arity counts $7,45,43,9$. Therefore a factorization with at least five distinct nonidentity residue classes modulo27 and a prime divisor $2\bmod3$ necessarily gives a hit in this shell. This is a pointwise sufficient consequence, not just a classification after a witness is given.

The bound nine is attained in the local profile with sole residue 23 and
capacity 9: its first E raw target exponent is 17 and its first M centered
target exponent has absolute value 9. This optimality is for the stated local
model, not a claim of prime optimality.

\paragraph{Marked inverse and arithmetic sharpness.}
Given $u$, put $d=\gcd(a,u)$ and $(h,r,s)=(d^2/u,u/d,a/d)$.
Prime valuations give $a=hrs$, $u=hr^2$, and $\gcd(r,s)=1$.
The inherited ordered formulas are
\[
 E:\ \kappa=(nr+s)/R,\ (a,hs\kappa,nhr\kappa),\qquad
 M:\ \lambda=(r+s)/R,\ (a,nhs\lambda,nhr\lambda).
 \tag{5}
\]
The target divisibility makes the quotients integral. Clearing reciprocals
gives $4/n$ since $n+R=4a$; the inverse is
$u=a^2/(RY-na)$ in E and $u=na^2/(RY-na)$ in M \cite[\S2]{ET}.
Thus the full denominator order, channel, scale and original prime exponents
remain available after capacity selection.

At the prime $p=19489\equiv169\pmod{840}$, $a=4879=7\cdot17\cdot41$,
$R=27$, every proper residue subbudget misses both targets. The full 54-state
shell has the sole E hit
\[
 u=580601,\quad(h,r,s,\kappa)=(41,119,1,85896),
 \quad(X,Y,Z)=(4879,3521736,8167578435576).
\]
This makes three classes necessary. Its E centered vector is $(1,1,0)$;
the arity theorem uses the raw E word $(2,2,1)$, not that alternative
coordinate count. Trial division through 139 certifies primality.

\paragraph{An infinite family with a hand proof.}
If all prime residues of $a$ lie in $\{1,7,11\}$, let $A,B$ be the total
multiplicities in classes 7 and 11. Then E is nonempty exactly when
$B\ge7$ or $(A\ge1,B\ge3)$ or $(A\ge2,B\ge1)$; M is nonempty exactly when
$B\ge9$ or $(A\ge1,B\ge5)$ or $(A\ge2,B\ge1)$.
Indeed $\log_2 7=-2$, $\log_2 11=-5$ modulo 18. For E solve
$2i+5j=11\pmod{18}$ with $i\in[0,2A]$, $j\in[0,2B]$:
$A=0$ needs $j=13$; $A=1$ first permits $j=5$; $A\ge2$ permits $(3,1)$;
$B=0$ is impossible. For M solve $2r+5s=9\pmod{18}$ in the centered box:
$A=0$ needs $|s|=9$; $A=1$ first permits $|s|=5$; $A\ge2$ permits $(2,1)$.
These cases prove both necessity and sufficiency for every multiplicity.

\paragraph{Scope.}
The complete scripts, 174-entry JSON catalogue, all inverse data and the
longer proof are supplied together. The application to 4,519 hard-class primes
through two million finds 1,214 occupied residual-27 shells, 1,880 E states and
1800 M states; it is not new prime coverage. The all-exponent theorem is at
fixed residual, not a proof that every prime has an occupied residual.
General group periodicity and positive rational series are classical tools.
The exact scalar comparison and finite catalogue were calculated here;
no historical-priority, Lean-build or independent-review claim is made.

\begin{thebibliography}{3}\small
\bibitem{SZ} Split-Zero research programme (2026).
\emph{Derived Mathematics}, \texttt{formal/splitzero/DERIVED\_MATHEMATICS.md},
\S\S1--4, revision \texttt{98c3c5f\ldots afffae}; and the owner's saved
\emph{Globalization Semiring Analysis}, \S6.1--6.2. Exact locators accompany the package.
\bibitem{Target} Clankers (2026). \emph{Support-Completed Center Sectors,
Koide Null Forms, and Fourier-Zero Removal in Residual Erd\H{o}s--Straus
Shells}. May source edition, Abstract; ordinary target $D_k=[3]+2[3-k]$.
\bibitem{ET} Elsholtz, C., \& Tao, T. (2013). Counting the number of solutions
to the Erd\H{o}s--Straus equation on unit fractions.
\emph{Journal of the Australian Mathematical Society, 94}(1), 50--105.
\S2, Propositions 2.2 and 2.6; arXiv:1107.1010.
\end{thebibliography}
\end{document}
